% Partie : histoire
% Chapitre : biographies
% Section : Alonzo Church

\documentclass[../../../include/open-logic-section]{subfiles}

\begin{document}

\olfileid[fr]{his}{bio}{chu}

\olsection{Alonzo Church}

\olphoto{church-alonzo}{Alonzo Church}

Alonzo Church naquit
à Washington,
dans le district
de Columbia,
le 14 juin 1903.
Dans sa petite enfance,
un accident avec
un pistolet à
air comprimé
le rendit aveugle
d'un œil.
Il termina ses
études préparatoires
dans le Connecticut
en 1920 et
commença la même
année ses études
universitaires
à Princeton.
Il acheva son
doctorat en 1927.
Après quelques
années à l'étranger,
il revint à
Princeton.
Church était
connu pour sa
politesse extrême
et sa minutie.
Son écriture
au tableau
était impeccable.
Pour conserver
des documents
importants,
il les recouvrait
soigneusement
de colle
transparente
Duco.
En dehors de
ses travaux,
il aimait lire
des magazines
de science-fiction
et n'hésitait
pas à écrire
à leurs rédacteurs
lorsqu'il repérait
une inexactitude
dans un texte.

Les contributions
scientifiques
de Church
furent considérables.
Avec ses étudiants
Stephen Kleene
et Barkley Rosser,
il élabora
une théorie
de la calculabilité
effective, le
lambda-calcul,
indépendamment
de la machine
de Turing
développée par
Alan Turing.
On démontre que ces deux modèles définissent les mêmes
fonctions. Selon la \emph{thèse de
Church--Turing}, une fonction sur les entiers naturels est
effectivement calculable, au sens intuitif, si et seulement
si elle est calculable par une machine de Turing ou, de façon
équivalente, au moyen du lambda-calcul. Church démontra
également le \emph{théorème de
Church}~: il n'existe pas de procédure effective qui décide
de la validité des formules du premier ordre.

Church poursuivit
ses recherches
jusqu'à un
âge avancé.
En 1967,
il quitta
Princeton pour
l'UCLA, où
il fut professeur
jusqu'à sa
retraite en 1990.
Il mourut
le 11 août 1995,
à l'âge
de 92 ans.

\begin{reading}
Pour une courte
biographie de Church,
voir \citet{EndertonND}.
Les articles originaux de
\citet{Church1936,Church1936a} portent sur le lambda-calcul
et le problème de la décision (\emph{Entscheidungsproblem}).
\citet{Aspray1984}
reproduit un entretien
avec Church sur
le milieu mathématique
de Princeton
dans les années 1930.
De 1936 à 1979,
Church rédigea
des comptes rendus
d'ouvrages pour
le \emph{Journal of
Symbolic Logic}.
Ils sont tous
archivés sur
le site de
John MacFarlane
\citep{MacFarlane2015}.
\end{reading}

\end{document}
